
\newpage
\section{Equipa, Gestão e Integração Contínua}

\subsection{Papéis da Equipa}
A equipa dividiu as responsabilidades e o foco de desenvolvimento da seguinte forma:
\begin{itemize}[nosep]
    \item \textbf{Brandon Mejia:} \textit{Fullstack}
    \item \textbf{Tiago Antunes:} \textit{Frontend}
    \item \textbf{Luís Moreira (Daniel):} \textit{Backend}
    \item \textbf{Miguel Correia:} \textit{Backend}
\end{itemize}

\subsection{Gestão de Projeto (SCRUM e Requisitos)}
O projeto segue a metodologia ágil SCRUM. Optámos por não duplicar documentação neste relatório, uma vez que o levantamento de requisitos, a definição de atores, a gestão do \textit{backlog} (com \textit{user stories}) e o planeamento dos \textit{sprints} estão totalmente centralizados e documentados no \textbf{Jira}.

Ao nível de código, utilizamos o Git com uma política de \textit{Pull Requests}, onde o código é revisto antes de ser integrado na \textit{branch} principal.

\subsection{Integração Contínua (CI)}
Implementámos \textit{pipelines} de CI via GitHub Actions para garantir a estabilidade constante do código. Foram criados dois ficheiros de \textit{workflow} distintos:
\begin{itemize}[nosep]
    \item \textbf{Backend CI}: Executa a construção e todos os testes unitários e de integração em Java.
    \item \textbf{Frontend CI}: Corre o \textit{linter} e os testes E2E/API usando Playwright.
\end{itemize}
Estes \textit{workflows} são disparados automaticamente a cada \textit{commit} e \textit{Pull Request}, prevenindo que código quebrado afete o ecossistema.